\SourceLeaf{382}%
\hspace*{\parindent}The more unmistakable the anomalous pattern here, the more interesting it is to ask whether the astronomically determined azimuths of the sides of the triangles show similar anomalies. In fact, according to a theorem that \textsc{Laplace} stated under particular restrictions, but that is capable of wide generalization, the convergence of the meridians has a necessary relation to the difference in longitude, independent of the figure of the Earth. so that, as one proceeds along a chain of geodetic lines, the irregularities of the one can be calculated \textit{a priori} from those of the other. As we have reported, the astronomical determinations of azimuth at the four principal stations, Milan, Turin, Mont Cenis and Colombier, were made with great care. The authors therefore compared the orientations at the last three stations with those obtained by transferring the Milan orientation by means of the geodetic measurements, using the same dimensions of the Earth's spheroid as before. The differences are
\begin{center}\begin{tabular}{lr}
for Turin & \(-5''5\)\\
Mont Cenis & \(-51,2\)\\
Colombier & \(-25,2\)
\end{tabular}\end{center}
Here too, then, one sees extraordinarily large anomalies. If, however, one calculates the anomalies in the differences in longitude from these figures by the theorem mentioned above (dividing by fifteen times the sine of the latitude of Milan and changing the sign), the resulting values are entirely different from those found directly, namely
\begin{center}\begin{tabular}{l|r|r}
& \multicolumn{1}{c|}{Calculated} & \multicolumn{1}{c}{Difference from the}\\
& \multicolumn{1}{c|}{anomaly} & \multicolumn{1}{c}{observed anomaly}\\
Turin & \(+0''52\) & \(+2''60\)\\
Mont Cenis & \(+4,81\) & \(+4,11\)\\
Colombier & \(+2,34\) & \(+1,95\)
\end{tabular}\end{center}
The authors remark of \textit{these} differences only that they are too large to be attributed to the accumulation of errors in the angle measurements, and thus leave us in the dark about what to make of them. In our view these three numbers are of the greatest importance,
